\section{Introduction}\label{sec:introduction}

Machine learning models have become essential components of modern cybersecurity infrastructure, enabling automated detection of intrusions, anomalies, and advanced persistent threats. However, the effectiveness and reliability of these models depend critically on the availability of rich, labeled, and realistic training datasets. In real-world security systems, malicious events are inherently rare, dynamic, and context-specific, creating substantial challenges for collecting sufficient training data to develop robust detection systems.

While several benchmark datasets exist, including CICIDS2017~\cite{sharafaldin2018toward} and UNSW-NB15~\cite{moustafa2015unsw}, they suffer from inherent limitations that constrain their utility for comprehensive model evaluation. These datasets are fundamentally static, lacking the configurability needed to systematically evaluate model robustness under varying threat levels or class imbalance scenarios that characterize real-world cybersecurity environments. Moreover, operational cybersecurity data is frequently inaccessible due to privacy constraints, legal restrictions, and confidentiality requirements, further exacerbating the data scarcity problem.

Synthetic data generation has emerged as a promising approach to address these challenges, with Large Language Models (LLMs) showing particular promise for generating realistic cybersecurity content. Recent advances in generative models have demonstrated substantial capability for producing high-fidelity synthetic data across various domains. Rahman et al.~\cite{rahman2025leveraging} explore GAN-based approaches for synthetic attack sample generation, while Razmyslovich et al.~\cite{razmyslovich2025eltex} introduce ELTEX, an LLM-based framework for domain-driven synthetic text generation in cybersecurity contexts.

However, current approaches to LLM-based cybersecurity synthetic data generation suffer from significant methodological limitations. Existing work typically relies on ad-hoc prompt design without systematic evaluation of prompt strategy effectiveness or comprehensive assessment of resulting data quality. Researchers often employ single prompt configurations based on intuition rather than empirical evidence, leading to inconsistent results and limited reproducibility. This gap between the demonstrated potential of LLM capabilities and the lack of systematic prompt engineering represents a critical bottleneck in scaling synthetic data solutions for cybersecurity applications.

The cybersecurity domain presents unique challenges for synthetic data generation that distinguish it from general text generation tasks. Cybersecurity data requires precise preservation of attack patterns, social engineering techniques, and domain-specific characteristics that determine malicious intent. Effective synthetic generation must balance maintaining semantic authenticity with introducing sufficient variation to improve model generalization capabilities. However, the relationship between prompt design choices and resulting synthetic data quality has not been systematically studied, leaving practitioners without evidence-based guidance for optimizing synthetic generation approaches.

Furthermore, the evaluation of synthetic cybersecurity data quality presents additional challenges beyond those encountered in other domains. Traditional metrics focused on lexical similarity or fluency may fail to capture the preservation of critical security-relevant characteristics that determine practical utility for classifier training. The complex interplay between synthetic data characteristics and classifier performance under realistic operational conditions, particularly under class imbalance scenarios common in cybersecurity applications, remains poorly understood.

To address these methodological gaps, this paper presents a systematic framework for evaluating LLM-generated synthetic data in cybersecurity applications. Our work investigates the fundamental research question: How do different prompt engineering strategies affect the quality and effectiveness of LLM-generated cybersecurity synthetic data for training robust detection models?

We examine this question through a comprehensive four-stage evaluation framework that progresses from systematic seed selection through controlled synthetic generation to rigorous performance assessment. Our approach enables systematic comparison of different prompt strategies while accounting for the unique characteristics and operational constraints of cybersecurity applications.

Specifically, we evaluate three distinct prompt engineering strategies across multiple dimensions of variation. The original strategy serves as a balanced baseline that preserves semantic content while enabling controlled variation. The strong strategy explores aggressive amplification of malicious characteristics, while the weak strategy investigates conservative approaches that maintain subtlety. This systematic comparison enables us to establish empirical relationships between prompt design choices and practical utility for cybersecurity classifier training.

Our experimental evaluation encompasses comprehensive analysis across 16 dataset configurations and three model architectures, providing insights into synthetic data effectiveness under varying operational conditions. Through systematic manipulation of class imbalance ratios from severe imbalance scenarios to more balanced conditions, we evaluate synthetic data utility across the spectrum of conditions encountered in real-world cybersecurity deployments.

This research makes several key contributions to the field of synthetic cybersecurity data generation:

\begin{enumerate}
\item \textbf{Systematic Evaluation Framework}: We present the first comprehensive framework for evaluating LLM-generated synthetic cybersecurity data that accounts for both intrinsic data quality and extrinsic classifier performance across multiple operational conditions.

\item \textbf{Empirical Prompt Strategy Analysis}: Through systematic comparison of three distinct prompt engineering approaches, we provide evidence-based guidance for optimizing LLM-based synthetic generation in cybersecurity contexts.

\item \textbf{Quality Assessment Methodology}: We establish embedding space analysis as an effective proxy for predicting synthetic data quality, enabling efficient optimization of generation strategies without expensive classifier training.

\item \textbf{Operational Insight}: Our evaluation under realistic class imbalance conditions provides practical guidance for deploying synthetic data augmentation in resource-constrained cybersecurity environments.
\end{enumerate}

Our findings demonstrate that synthetic data can achieve F1-scores within 0.012 of real data performance under optimal conditions, with particular value in low-resource scenarios where real malicious samples are scarce. The systematic framework we present establishes benchmarks for evaluating synthetic cybersecurity data while providing standardized protocols that enable meaningful comparison across different approaches.

The remainder of this paper presents our four-stage evaluation framework, systematic experimental design, comprehensive results analysis, and implications for future research in synthetic cybersecurity data generation. All experimental protocols and evaluation metrics are designed to support reproducibility and enable cumulative progress in this rapidly evolving field.